% Cover option L: "rotation curve" -- dusk style; a spiral galaxy's flat
% rotation curve vs. the falling Keplerian prediction, the gap labeled as
% dark matter.
\documentclass[11pt,a4paper]{article}
\usepackage[margin=0cm]{geometry}
\usepackage[T1]{fontenc}
\usepackage{lmodern}
\usepackage{tikz}
\usetikzlibrary{calc}
\pagestyle{empty}

\definecolor{gtop}{HTML}{170D2B}
\definecolor{gbottom}{HTML}{51285F}
\definecolor{gpink}{HTML}{F48FC0}
\definecolor{gmagenta}{HTML}{E0559E}
\definecolor{glav}{HTML}{CBB4DD}
\definecolor{gwhite}{HTML}{F8F3FB}

\begin{document}
\begin{tikzpicture}[remember picture,overlay]
  \shade[top color=gtop,bottom color=gbottom]
    (current page.south west) rectangle (current page.north east);

  \begin{scope}[shift={(current page.south west)}]
    \begin{scope}[shift={(2.8,5.2)}]
      % Faint axes.
      \draw[glav,opacity=0.35,line width=0.5pt,
            dash pattern=on 1.6pt off 2.2pt,-stealth] (-0.3,0) -- (16.4,0);
      \draw[glav,opacity=0.35,line width=0.5pt,
            dash pattern=on 1.6pt off 2.2pt,-stealth] (0,-0.3) -- (0,8.4);
      \node[glav,opacity=0.7,font=\footnotesize\sffamily]
        at (14.4,-0.55) {distance from centre};
      \node[glav,opacity=0.7,font=\footnotesize\sffamily,rotate=90]
        at (-0.5,6.6) {orbital speed};
      % A little edge-on galaxy at the origin.
      \fill[glav,opacity=0.6] (0.45,0.5) ellipse (0.62 and 0.13);
      \fill[gwhite] (0.45,0.5) circle (0.09);
      % Keplerian expectation from the visible matter, falling away.
      \draw[glav,opacity=0.6,line width=0.8pt,smooth,
            dash pattern=on 2.4pt off 2.6pt]
        plot[domain=3.0:15.6,samples=90] (\x,{9.35/sqrt(\x)});
      \node[glav,font=\footnotesize\sffamily,anchor=west] at (11.9,1.95)
        {visible matter alone};
      % The observed curve: climbs, then stays flat (with mild wiggles).
      \draw[gpink,line width=1.3pt,smooth]
        plot[domain=0.02:15.6,samples=140]
        (\x,{6.3*\x/sqrt(\x*\x+3.24) + 0.1*sin(\x*80)});
      % The gap is the point of the plot.
      \draw[gmagenta,line width=0.8pt,stealth-stealth]
        (13.6,{9.35/sqrt(13.6)+0.18}) -- (13.6,{6.3*13.6/sqrt(13.6*13.6+3.24)
          + 0.1*sin(13.6*80) - 0.18});
      \node[gmagenta,font=\footnotesize\sffamily,anchor=west] at (13.85,4.35)
        {dark matter};
      % Observed marker, glowing; a white dot on the Keplerian curve.
      \pgfmathsetmacro{\vo}{6.3*11.0/sqrt(11.0*11.0+3.24)+0.1*sin(11.0*80)}
      \fill[gmagenta,opacity=0.22] (11.0,\vo) circle (0.62);
      \fill[gmagenta,opacity=0.4] (11.0,\vo) circle (0.4);
      \fill[gpink] (11.0,\vo) circle (0.22);
      \node[gpink,font=\footnotesize\sffamily] at (11.0,\vo+0.75) {observed};
      \fill[gwhite] (8.0,{9.35/sqrt(8.0)}) circle (2.6pt);
    \end{scope}
  \end{scope}

  % ---------------- title block, left ----------------
  \begin{scope}[shift={(current page.north west)}]
    \node[anchor=west,gpink,font=\small\sffamily\scshape] at (2.2,-3.1)
      {International Olympiad on Astronomy \& Astrophysics};
    \draw[gmagenta,line width=1.6pt] (2.24,-3.75) -- (7.0,-3.75);
    \node[anchor=west,gwhite,font=\fontsize{40}{44}\selectfont\sffamily\bfseries]
      at (2.15,-5.4) {IOAA};
    \node[anchor=west,gwhite,font=\fontsize{40}{44}\selectfont\sffamily\bfseries]
      at (2.15,-7.0) {Problem Book};
    \node[anchor=west,glav,font=\large\sffamily,align=left] at (2.2,-9.4)
      {Problems and available official solutions\\ organized by topic};
    \node[anchor=west,gpink,font=\LARGE\sffamily\bfseries] at (2.2,-11.4)
      {2007--2025};
  \end{scope}
  \node[anchor=west,glav,font=\normalsize\sffamily]
    at ($(current page.south west)+(2.2,2.0)$)
    {Compiled and edited by Ioana A.~Zelko};
\end{tikzpicture}
\end{document}
